\documentclass[11pt]{article}
\usepackage[a4paper,margin=0.82in]{geometry}
\usepackage{amsmath,amssymb,amsthm,mathtools}
\usepackage{booktabs,longtable,tabularx,array}
\usepackage{graphicx}
\usepackage{xcolor}
\usepackage{hyperref}
\usepackage{enumitem}
\usepackage{caption}
\usepackage{subcaption}
\usepackage{fancyhdr}
\usepackage{microtype}
\usepackage{titlesec}
\usepackage{lastpage}
\usepackage{listings}

\definecolor{accent}{HTML}{174A7C}
\definecolor{soft}{HTML}{EEF5FB}
\definecolor{dark}{HTML}{1F2937}
\hypersetup{colorlinks=true,linkcolor=accent,citecolor=accent,urlcolor=accent}
\graphicspath{{figures/}}
\titleformat{\section}{\Large\bfseries\color{accent}}{\thesection}{0.7em}{}
\titleformat{\subsection}{\large\bfseries\color{dark}}{\thesubsection}{0.7em}{}
\pagestyle{fancy}
\fancyhf{}
\lhead{Referral auctions on networks}
\rhead{Final project}
\cfoot{\thepage/\pageref{LastPage}}

\newtheorem{definition}{Definition}
\newtheorem{theorem}{Theorem}
\newtheorem{lemma}{Lemma}
\newtheorem{proposition}{Proposition}
\newtheorem{corollary}{Corollary}
\newtheorem{remark}{Remark}
\newtheorem{conjecture}{Conjecture}

\newcommand{\R}{\mathbb{R}}
\newcommand{\E}{\mathbb{E}}
\newcommand{\N}{\mathcal{N}}
\newcommand{\A}{\mathcal{A}}
\newcommand{\Rev}{\operatorname{Rev}}
\newcommand{\SW}{\operatorname{SW}}
\newcommand{\argmax}{\operatorname*{arg\,max}}
\newcommand{\score}{\operatorname{score}}
\newcommand{\PathOp}{\operatorname{path}}
\newcommand{\subtree}{\operatorname{subtree}}

\begin{document}

\begin{titlepage}
    \centering
    \vspace*{1.0cm}
    {\Huge\bfseries Incentive-Compatible Referral Auctions on Networks\\[0.25cm]}
    {\Large Cost, Sybil, Divisible, Multi-Item, and Multi-Seller Extensions\\[0.8cm]}
    \begin{tabular}{rl}
    Team: & Abhi Jain (23B0903), Harshit Raj (23B0950), Ridam Jan (23B1049)\\
    Baseline: & Li, Hao, Zhao, Zhou, \emph{Mechanism Design in Social Networks}\\
    Deliverable: & Theory report + runnable simulator + plots\\
    \end{tabular}
    \vfill
    \fcolorbox{accent}{soft}{\begin{minipage}{0.90\textwidth}
    \textbf{Final positioning.} The project is not framed as ``we solved all diffusion auctions''.  We use the baseline IDM paper as the core model, identify realistic variants, separate known literature from open-ish gaps, prove one clean extension in a restricted domain, and implement stress tests for the broader variants.  The main theorem-level contribution is a participation-cost version of IDM; the multi-item, divisible, sybil, and multi-seller parts are modeled carefully and evaluated experimentally, with honest boundaries.
    \end{minipage}}
    \vfill
    {\large April 2026}
\end{titlepage}

\tableofcontents
\newpage

\section{Executive summary}

Classical auctions assume that the seller can directly reach all bidders.  The baseline paper of Li et al. introduced a different setting: the seller initially informs only her neighbors in a social graph, and each informed buyer chooses both a bid and a set of neighbors to invite.  The hard incentive problem is that inviting neighbors may create competitors.  Network-VCG can incentivize diffusion but may run a deficit; IDM restores weak budget balance by paying only diffusion-critical agents in a carefully thresholded way~\cite{li2017mdsn,li2022diffusion}.

The requested variants are broad.  Some are already partially addressed by later literature: multiple items, multi-unit auctions, combinatorial diffusion auctions, transaction costs, sybil attacks, budgeted propagation, and optimal referral auctions all have recent papers~\cite{zhao2018multi,fang2023multiunit,li2024transaction,chen2023sybil,zheng2024budget,bhattacharyya2024optimal,fang2024meta,li2024combinatorial}.  A good course project should therefore not merely reimplement known mechanisms.  We propose the following final alignment:

\begin{center}
\begin{tabularx}{0.96\textwidth}{p{0.24\textwidth}X p{0.24\textwidth}}
\toprule
\textbf{Scenario} & \textbf{What we do} & \textbf{Claim status}\\
\midrule
Participation costs & Develop Participation-Cost IDM (PC-IDM), prove IC/IR/WBB by reducing to IDM on net values. & Main theorem contribution\\
Sybil attacks & Stress-test VCG/IDM under fake identities; implement tax baseline; cite STM/SCM as stronger known mechanisms. & Simulation + conditional lemma\\
Continuous divisible item & Model as limit of many homogeneous units; prove VCG deficit persists; simulate discrete approximation. & Model + negative result + simulator\\
Multiple different items & Use public-quality heterogeneous item domain; implement VCG benchmark; connect to multi-item/combinatorial literature. & Restricted benchmark + future work\\
Multiple sellers & Fixed partition into seller regions and run PC-IDM locally. & Conditional theorem + simulation\\
\bottomrule
\end{tabularx}
\end{center}

\section{Baseline model and notation}

Let $G=(V,E)$ be a graph with seller $s\in V$ and buyers $N=V\setminus\{s\}$.  Buyer $i$ has valuation $v_i\ge 0$ for the item.  Initially only the seller's neighbors know the auction.  Buyer $i$ reports
\[
    a_i=(b_i,r_i'),\qquad r_i'\subseteq r_i,
\]
where $b_i$ is the value report and $r_i'$ is the set of neighbors invited.  A feasible action profile contains only buyers reachable from $s$ in the induced invitation graph.

A mechanism $M=(\pi,p)$ chooses allocation $\pi_i(a)\in\{0,1\}$ and payments $p_i(a)\in\R$, where positive $p_i$ means $i$ pays the seller and negative $p_i$ means $i$ is rewarded.  Buyer utility is quasilinear:
\[
    u_i = \pi_i(a)v_i - p_i(a).
\]

\begin{definition}[Diffusion incentive compatibility]
A mechanism is diffusion incentive-compatible if, for every buyer, truthfully reporting the value and inviting all neighbors is a dominant strategy.
\end{definition}

\begin{definition}[Weak budget balance]
A mechanism is weakly budget balanced if $\Rev(M,a)=\sum_i p_i(a)\ge 0$ for every feasible profile $a$.
\end{definition}

For a rooted tree, buyer $i$ is a diffusion-critical node of buyer $j$ exactly when $i$ lies on the unique path from $s$ to $j$.  Hence
\[
    d_i=\{i\}\cup\{j: i\text{ is diffusion-critical for }j\}=\subtree(i).
\]
Let $v^*_{-d_i}$ denote the highest bid among buyers outside $d_i$; it is $0$ if the set is empty.

\subsection{Known mechanisms used as baselines}

\paragraph{Local Vickrey.} The seller runs a second-price auction only among direct neighbors.  It is safe but ignores deeper high-value buyers.

\paragraph{Network VCG.} Allocate to the highest reachable bidder and charge each buyer the externality caused by her presence:
\[
    p_i^{VCG}=W(a_{-d_i})-\left(W(a)-\pi_i(a)b_i\right).
\]
It is truthful for value and diffusion, but may have negative seller revenue on a path.

\paragraph{IDM.} Let $m$ be the highest bidder and let $C_m=(c_1,c_2,\ldots,m)$ be the diffusion-critical sequence of $m$.  IDM allocates to the first node on this path that would be the highest bidder if the next critical block were removed; otherwise it allocates to $m$.  If $w$ wins, then for $i\in C_w\setminus\{w\}$,
\[
    p_i^{IDM}=v^*_{-d_i}-v^*_{-d_{i+1}},
\]
while the winner pays
\[
    p_w^{IDM}=v^*_{-d_w}.
\]
The on-path payments are nonpositive rewards.  IDM is IC, IR, and WBB, and its revenue is at least network-VCG revenue in the baseline model~\cite{li2017mdsn}.

\section{Literature map and project gaps}

\begin{longtable}{p{0.19\textwidth}p{0.31\textwidth}p{0.22\textwidth}p{0.21\textwidth}}
\toprule
\textbf{Paper / direction} & \textbf{Known result} & \textbf{Assumptions} & \textbf{Gap useful for us}\\
\midrule
Li et al. IDM~\cite{li2017mdsn,li2022diffusion} & First social-network auction model, network VCG, IDM, critical diffusion ideas. & Single item; no participation cost; no sybil identities. & Add realistic frictions while preserving a reduction to IDM.\\
\addlinespace
Selling multiple items~\cite{zhao2018multi} & Generalizes diffusion incentives to multiple homogeneous items. & Multiple identical items. & Heterogeneous items and divisible limits need separate modeling.\\
\addlinespace
Multi-unit/MUDAN~\cite{fang2023multiunit} & Multi-unit diffusion with graph exploration, IC and weak efficiency guarantees. & Homogeneous items, multi-demand. & Use as comparison for discretized divisible allocation.\\
\addlinespace
Combinatorial diffusion~\cite{fang2024meta,li2024combinatorial} & Meta-mechanisms/DCAF for combinatorial domains; preserve IC/IR/WBB under assumptions. & Single-minded or monotone valuation models, framework assumptions. & Full arbitrary externalities remain too broad; restrict to public-quality item values.\\
\addlinespace
Sybil-proof diffusion~\cite{chen2023sybil} & STM and SCM address fake identities in single-item diffusion auctions. & Single item, sybil attack model. & We can stress-test baseline mechanisms and implement a simple tax proxy.\\
\addlinespace
Transaction costs~\cite{li2024transaction} & Diffusion auctions with transaction costs and critical neighborhood ideas. & Intermediaries and transaction-cost framework. & Distinguish participation cost from forwarding cost; prove an impossibility for full costly diffusion.\\
\addlinespace
Optimal referral auctions~\cite{bhattacharyya2024optimal} & Myerson-style characterization; optimal referral auctions under i.i.d.; experiments for non-i.i.d. & Single item, referral auction class. & Non-i.i.d. and topology robustness are good experimental axes.\\
\bottomrule
\end{longtable}

\section{Main extension: participation-cost IDM}

\subsection{Model}

In many real auctions, bidders can tolerate a small participation cost: time, platform fees, verification, or effort.  Let $c_i\ge 0$ be a public participation cost for buyer $i$.  The utility is
\[
    u_i=\pi_i v_i-p_i-c_i\cdot \mathbf{1}\{i\text{ participates}\}.
\]
We assume $c_i$ is incurred once after participation and is not a per-neighbor forwarding cost.  This distinction is important: per-edge forwarding costs create an impossibility discussed later.

Define the net value
\[
    y_i=(v_i-c_i)^+=\max\{v_i-c_i,0\}.
\]

\begin{definition}[Participation-Cost IDM]
PC-IDM runs IDM on the transformed single-parameter reports
\[
    \hat y_i=(b_i-c_i)^+.
\]
Allocation is computed exactly as IDM using $\hat y_i$.  Payments are the IDM threshold payments in the transformed score space.
\end{definition}

\subsection{Properties}

\begin{lemma}[Score reduction]
Fix public costs $c_i$.  For every buyer $i$, reporting $v_i$ truthfully is equivalent to reporting the transformed net value $y_i=(v_i-c_i)^+$ truthfully.
\end{lemma}

\begin{proof}
The map $v_i\mapsto (v_i-c_i)^+$ is monotone and depends only on the public constant $c_i$.  Therefore any profitable misreport in $v_i$ would induce a profitable misreport in the single-parameter score $y_i$, and conversely any score misreport can be implemented by some value report $b_i=y_i+c_i$ when $y_i>0$.
\end{proof}

\begin{theorem}[PC-IDM is truthful and weakly budget balanced]
On the baseline social-network auction model with public participation costs, PC-IDM is value incentive-compatible, diffusion incentive-compatible, individually rational with respect to net utility, and weakly budget balanced.
\end{theorem}

\begin{proof}
Run IDM on the transformed nonnegative scores $y_i$.  The baseline IDM theorem applies to any nonnegative single-parameter valuations: truthful score reporting and full diffusion are dominant, every truthful participant gets nonnegative score-utility, and total payment is nonnegative.  By the score reduction lemma, truthful reporting of $v_i$ is equivalent to truthful reporting of $y_i$.  Since the actual utility equals score utility,
\[
    \pi_i v_i-p_i-c_i = \pi_i(v_i-c_i)-p_i
\]
for winners and equals the corresponding IDM reward utility for non-winners, truthful utility is nonnegative whenever the transformed allocation gives positive score.  If $v_i<c_i$, the transformed value is zero, so PC-IDM never requires that buyer to win at a loss.  Finally, the sum of PC-IDM payments is the sum of IDM payments on $y$, hence nonnegative.
\end{proof}

\begin{remark}
This is a modest but defensible extension: we do not solve arbitrary costly diffusion.  We show that a common real-world friction can be absorbed into the IDM threshold structure without breaking IC or WBB.
\end{remark}

\subsection{Why per-edge forwarding cost is harder}

\begin{proposition}[Impossibility for full diffusion with positive edge costs]
Suppose every invited edge costs the inviter $\kappa>0$, and the mechanism must be ex-post IR, ex-post WBB, and make inviting all neighbors a dominant strategy for every valuation profile.  No such mechanism exists without an outside subsidy.
\end{proposition}

\begin{proof}
Consider any graph where a buyer has at least one neighbor to invite, and consider the valuation profile $v_i=0$ for all buyers.  Full diffusion requires that the buyer prefers paying the edge cost and inviting the neighbor.  Ex-post IR then requires the mechanism to compensate at least $\kappa$ for this action.  Since all valuations are zero, no buyer can be charged a positive amount while preserving ex-post IR.  Thus the seller must pay a positive reward while receiving no positive payment, violating WBB.
\end{proof}

This is why the simulator separates participation costs, where a clean theorem is possible, from transaction/forwarding costs, where the right mechanism must use reserves, budgets, or weaker diffusion requirements.

\section{Scenario: continuous divisible allocation}

\subsection{Model}

The seller has a divisible mass $Q=1$.  Buyer $i$ can receive $x_i\in[0,\bar x_i]$ and has linear utility $v_i x_i$ before transfers.  The feasible set is
\[
    \sum_i x_i\le 1,\qquad 0\le x_i\le \bar x_i.
\]
A full continuous diffusion auction should choose $x$ and payments while incentivizing both truthful value reports and information diffusion.

\subsection{What is known and what is open}

Multi-unit diffusion auctions and fixed-price tree mechanisms exist for multiple homogeneous items~\cite{zhao2018multi,fang2023multiunit,yu2024tree}.  Divisible-good auctions are classical outside networks, but continuous divisible allocation with social-network diffusion is not simply IDM with fractions, because critical ancestors may enable many small allocations simultaneously.

\begin{proposition}[VCG deficit persists in the divisible limit]
The network-VCG deficit example from the baseline paper persists for a divisible item with linear values.
\end{proposition}

\begin{proof}
Use a path $s-1-2-\cdots-\ell$.  Let $v_\ell=1$ and all other values be zero.  The efficient divisible allocation gives the entire mass to $\ell$ and welfare $1$.  Every intermediate node is diffusion-critical for $\ell$.  Removing any intermediate node removes the only positive-value participant, so VCG pays that intermediate node a reward of $1$.  The winner pays zero.  Seller revenue is $-(\ell-1)$, as in the indivisible case.
\end{proof}

\subsection{Simulator treatment}

The code uses a $K$-micro-unit approximation.  For linear values and no demand caps, the efficient allocation of all micro-units goes to the highest bidder, so the divisible benchmark converges immediately.  With demand caps, this can be extended by allocating units greedily by marginal value.  Our conclusion is that continuous divisibility is a good future extension, but not the best theorem target unless we introduce demand caps or concave utilities.

\section{Scenario: multiple different items}

\subsection{Model}

Let there be heterogeneous items $M=\{1,\dots,m\}$.  General valuations $v_i(S)$ are combinatorial and can encode substitutes and complements.  This is now close to recent combinatorial diffusion-auction work~\cite{fang2024meta,li2024combinatorial}.  For a tractable project contribution, we study a restricted public-quality domain:
\[
    v_i(j)=q_j\theta_i,\qquad q_j>0,
\]
where each buyer wants at most one item, $q_j$ is public quality, and $\theta_i$ is a private scalar type.

\begin{proposition}[Truthfulness of the reached-set heterogeneous benchmark]
In the public-quality unit-demand domain, VCG over the reached set is dominant-strategy truthful in the scalar value reports.
\end{proposition}

\begin{proof}
The domain is a standard quasilinear private-value allocation problem once the reached set is fixed.  VCG maximizes reported welfare and charges the externality imposed on other reached buyers.  The standard Groves argument gives dominant-strategy truthfulness.  The restriction $v_i(j)=q_j\theta_i$ is single-parameter, but the claim follows even from the more general VCG theorem on the reached set.
\end{proof}

\begin{remark}
This is not a full diffusion mechanism for heterogeneous items.  It is a benchmark.  A full project extension would combine this restricted item model with a MetaMSN/DCAF-style transformation, then check whether non-sensitivity assumptions are satisfied.
\end{remark}

\section{Scenario: Sybil attacks}

A sybil attack means one strategic participant creates fake identities.  Recent work proposes Sybil Tax Mechanism and Sybil Cluster Mechanism for diffusion auctions~\cite{chen2023sybil}.  We do not rederive those mechanisms.  Instead, our project includes a stress test showing why the issue matters.

\subsection{Conditional tax lemma}

\begin{lemma}[Bounded sybil-gain tax]
Suppose inserting one fake identity can increase a strategic cluster's total transfer by at most $\Delta$ in a given mechanism class.  If every identity must pay a registration tax $\tau\ge \Delta$, then inserting fake identities is weakly unprofitable.
\end{lemma}

\begin{proof}
For $k$ fake identities, the total transfer gain is at most $k\Delta$ by the assumed marginal bound.  The total tax is $k\tau\ge k\Delta$.  Hence the net gain is nonpositive.
\end{proof}

This is not a replacement for STM/SCM: the hard part is proving a small $\Delta$ for a rich graph class.  The simulator therefore treats tax-IDM/tax-VCG as a baseline, not as a final sybil-proof mechanism.

\section{Scenario: multiple sellers}

Let $S\subset V$ be a set of sellers.  Each seller $s_r\in S$ owns one item.  Buyers may be reachable from multiple sellers.  A fully strategic model must handle seller competition, buyer choice, and possibly double-auction effects.  There is literature on social-network double auctions and related transaction-cost auctions~\cite{xu2020double,li2024transaction}.

For a feasible course extension, we use a fixed partition model: a public rule assigns each buyer to a seller region, and each seller runs PC-IDM inside its region.

\begin{proposition}[Fixed-region inheritance]
If the seller-region partition is public and independent of bids and invitations inside each region, then running PC-IDM independently in every region is IC, net-IR, and WBB within each region.
\end{proposition}

\begin{proof}
Condition on any fixed region.  The mechanism inside the region is exactly PC-IDM, so the previous theorem applies.  Payments and allocations in other regions are independent of a buyer's report in this region.  Summing nonnegative regional revenues gives global WBB.
\end{proof}

\begin{remark}
The assumption that the partition is independent of strategic reports is strong.  Relaxing it is a genuine future-work direction.
\end{remark}

\section{Experiments}

\subsection{Simulator design}

The simulator uses rooted trees, non-i.i.d. value distributions, public participation costs, sybil insertion, and modular mechanisms.  The mechanisms implemented are:
\begin{itemize}[leftmargin=*]
    \item Local Vickrey among seller-neighbors only;
    \item Network VCG with diffusion-critical payments;
    \item IDM on raw values;
    \item PC-IDM on net values $(v_i-c_i)^+$;
    \item tax variants for sybil stress tests;
    \item multi-unit VCG and public-quality heterogeneous VCG benchmarks;
    \item fixed-partition multi-seller PC-IDM.
\end{itemize}

\subsection{Numerical setup}

We run Monte-Carlo simulations with $n\in\{20,40,60\}$ buyers, random/path/binary rooted trees, non-i.i.d. Gaussian value classes, and gamma-distributed participation costs.  Each plotted point is an average over 20 independent runs.  The code also includes an exhaustive brute-force value-report check for PC-IDM on a toy path instance.

\subsection{Results}

\begin{figure}[h]
\centering
\includegraphics[width=0.86\textwidth]{single_item_cost_revenue_mean.pdf}
\caption{Seller revenue under participation costs.  PC-IDM preserves the IDM-style threshold structure but uses net values.  Local Vickrey loses deeper bidders; Network-VCG can have lower revenue because of critical-node rewards.}
\end{figure}

\begin{figure}[h]
\centering
\includegraphics[width=0.86\textwidth]{single_item_cost_net_welfare_mean.pdf}
\caption{Net welfare under participation costs.  PC-IDM is close to IDM in this sample because participation costs are moderate; Local Vickrey loses welfare by staying local.}
\end{figure}

\begin{center}
\begin{tabular}{lrr}
\toprule
\textbf{Mechanism, random topology, $n=60$} & \textbf{Revenue mean} & \textbf{Net welfare mean}\\
\midrule
Local Vickrey & 72.92 & 86.04\\
Network VCG & 98.40 & 113.55\\
IDM & 109.12 & 113.55\\
PC-IDM & 104.28 & 113.76\\
Tax-IDM & 169.12 & 113.55\\
\bottomrule
\end{tabular}
\end{center}

Tax-IDM is included only as a stress-test baseline: charging every identity increases revenue mechanically but changes participation incentives and is not claimed as a polished mechanism.

\begin{figure}[h]
\centering
\includegraphics[width=0.86\textwidth]{sybil_revenue.pdf}
\caption{Sybil stress test on a path.  Network-VCG becomes more negative as fake intermediates are inserted.  A per-identity tax shifts revenue upward but does not solve the full sybil-proofness problem.}
\end{figure}

\begin{figure}[h]
\centering
\begin{subfigure}{0.48\textwidth}
\includegraphics[width=\textwidth]{multi_item_revenue_mean.pdf}
\caption{Revenue}
\end{subfigure}
\hfill
\begin{subfigure}{0.48\textwidth}
\includegraphics[width=\textwidth]{multi_item_welfare_mean.pdf}
\caption{Welfare}
\end{subfigure}
\caption{Multi-unit and public-quality heterogeneous-item benchmarks.  Homogeneous multi-unit VCG has higher welfare here because the public-quality items have qualities $1,0.8,0.55$, so the heterogeneous benchmark deliberately has lower total value.}
\end{figure}

\begin{center}
\begin{tabular}{lrr}
\toprule
\textbf{Multi-item benchmark, $n=60$} & \textbf{Revenue mean} & \textbf{Welfare mean}\\
\midrule
Heterogeneous public-quality VCG & 249.81 & 270.90\\
Multi-unit VCG, $K=3$ & 314.57 & 342.29\\
\bottomrule
\end{tabular}
\end{center}

For the fixed-partition multi-seller experiment, aggregate revenue increases with network size: 146.08 at $n=20$, 170.40 at $n=40$, and 178.99 at $n=60$.  This is not a theorem about strategic seller competition; it only shows that the regional PC-IDM simulator is working and can be extended.

\section{What is novel in this final project?}

\begin{enumerate}[leftmargin=*]
    \item \textbf{A clean participation-cost extension.}  PC-IDM is not just a heuristic: it has a proof by reduction to IDM on transformed net values.
    \item \textbf{A useful impossibility boundary.}  Per-edge forwarding costs cannot be handled by demanding full diffusion, ex-post IR, and WBB for all value profiles without subsidy.
    \item \textbf{Careful separation of solved and unsolved regions.}  Multi-item, combinatorial, sybil-proof, and transaction-cost variants have existing papers.  We position our work as restricted positive theory plus simulation stress tests, not exaggerated claims.
    \item \textbf{A modular simulator.}  The code lets us plug in topologies, costs, sybil chains, item qualities, and seller regions.
\end{enumerate}

\section{Conclusions by scenario}

\begin{center}
\begin{tabularx}{0.96\textwidth}{p{0.22\textwidth}X}
\toprule
\textbf{Scenario} & \textbf{Conclusion}\\
\midrule
Participation cost & Best final-project target.  The model is realistic, the mechanism is simple, and the proof is clean.\\
Divisible allocation & Useful future work.  The VCG deficit survives, and discretization gives a simulator, but a nontrivial theorem needs demand caps or concave utilities.\\
Multiple different items & General case overlaps with combinatorial diffusion auction papers.  Our public-quality model is a defensible restricted benchmark.\\
Sybil attacks & Existing STM/SCM literature is directly relevant.  Our contribution is a stress-test and conditional tax lemma, not a complete sybil-proof mechanism.\\
Multiple sellers & Fixed-region PC-IDM inherits properties.  Strategic seller-region formation remains open and interesting.\\
\bottomrule
\end{tabularx}
\end{center}

\section{How to present this to the professor}

The strongest pitch is:
\begin{quote}
We began with the IDM baseline.  After the literature survey, we found that several natural extensions already have partial solutions: multi-item diffusion, combinatorial diffusion, sybil-proof diffusion, transaction costs, and optimal referral auctions.  Therefore we chose one narrow extension where we can prove something: public participation costs.  We prove PC-IDM by reduction to IDM, show an impossibility for full diffusion with per-edge costs, and implement simulations for the other real-world variants to identify future directions.
\end{quote}

This matches a research workflow: literature map, narrow gap, theorem, negative result, simulator, and honest future work.

\newpage
\begin{thebibliography}{99}

\bibitem{li2017mdsn}
Bin Li, Dong Hao, Dengji Zhao, and Tao Zhou.
\newblock Mechanism Design in Social Networks.
\newblock \emph{AAAI}, 2017. \url{https://arxiv.org/abs/1702.03627}

\bibitem{li2022diffusion}
Bin Li, Dong Hao, Hui Gao, and Dengji Zhao.
\newblock Diffusion auction design.
\newblock \emph{Artificial Intelligence}, 303:103631, 2022.

\bibitem{zhao2018multi}
Dengji Zhao, Bin Li, Junping Xu, Dong Hao, and Nicholas R. Jennings.
\newblock Selling Multiple Items via Social Networks.
\newblock \emph{AAMAS}, 2018 / arXiv version. \url{https://arxiv.org/abs/1903.02703}

\bibitem{fang2023multiunit}
Yuan Fang, Mengxiao Zhang, Jiamou Liu, Bakh Khoussainov, and Mingyu Xiao.
\newblock Multi-unit Auction over a Social Network.
\newblock arXiv:2302.08924, 2023. \url{https://arxiv.org/abs/2302.08924}

\bibitem{chen2023sybil}
Hongyin Chen, Xiaotie Deng, Ying Wang, Yue Wu, and Dengji Zhao.
\newblock Sybil-Proof Diffusion Auction in Social Networks.
\newblock \emph{AAMAS}, 2023. \url{https://arxiv.org/abs/2211.01984}

\bibitem{bhattacharyya2024optimal}
Rangeet Bhattacharyya, Parvik Dave, Palash Dey, and Swaprava Nath.
\newblock Optimal Referral Auction Design.
\newblock \emph{AAMAS}, 2024. \url{https://arxiv.org/abs/2208.09326}

\bibitem{fang2024meta}
Yuan Fang, Mengxiao Zhang, Jiamou Liu, and Bakh Khoussainov.
\newblock Meta-mechanisms for Combinatorial Auctions over Social Networks.
\newblock arXiv:2408.04555, 2024. \url{https://arxiv.org/abs/2408.04555}

\bibitem{li2024combinatorial}
Xuanyu Li, Miao Li, Yuhan Cao, and Dengji Zhao.
\newblock Combinatorial Diffusion Auction Design.
\newblock arXiv:2410.22765, 2024. \url{https://arxiv.org/abs/2410.22765}

\bibitem{li2024transaction}
Bin Li, Dong Hao, and Dengji Zhao.
\newblock Diffusion auction design with transaction costs.
\newblock \emph{Autonomous Agents and Multi-Agent Systems}, 2024. \url{https://dl.acm.org/doi/abs/10.1007/s10458-023-09631-8}

\bibitem{zheng2024budget}
Junjie Zheng, Xu Ge, Bin Li, and Dengji Zhao.
\newblock Sybil-Proof Mechanism for Information Propagation with Budgets.
\newblock arXiv:2405.14293, 2024. \url{https://arxiv.org/abs/2405.14293}

\bibitem{yu2024tree}
Feiyang Yu.
\newblock Diffusion Mechanism Design in Tree-Structured Social Network.
\newblock arXiv:2407.21143, 2024. \url{https://arxiv.org/abs/2407.21143}

\bibitem{xu2020double}
Junping Xu et al.
\newblock Design of Double Auction Mechanism Based on Social Network.
\newblock 2020.

\end{thebibliography}

\end{document}
